% figures/w05-switching.svg — 전환 판정 둘은 같은 판정이 아니다.
%
%   bash _tools/tikz2svg.sh figures/src/w05-switching.tex figures/w05-switching.svg
%
% 같은 배, 같은 자리. 남은 along-track 거리는 R 보다 작고(전환), 참 거리는
% R 의 두 배가 넘는다(대기). section F 가 이 숫자를 찍는다:
%   d - x_e^{\,p} = 4.50 m < R = 5,  |P - wp| = 10.97 m > R
\documentclass[border=5pt]{standalone}
% gnc-style.tex — 이 볼트의 TikZ 그림이 공유하는 스타일.
%
% 저널 그림 규약을 스타일 하나로 굳혀 둔다. 그림마다 선 굵기나 화살촉을
% 다시 정하지 않는다 — 그것이 그림이 제각각으로 보이는 원인이다.
%
% 쓰는 법:  % gnc-style.tex — 이 볼트의 TikZ 그림이 공유하는 스타일.
%
% 저널 그림 규약을 스타일 하나로 굳혀 둔다. 그림마다 선 굵기나 화살촉을
% 다시 정하지 않는다 — 그것이 그림이 제각각으로 보이는 원인이다.
%
% 쓰는 법:  % gnc-style.tex — 이 볼트의 TikZ 그림이 공유하는 스타일.
%
% 저널 그림 규약을 스타일 하나로 굳혀 둔다. 그림마다 선 굵기나 화살촉을
% 다시 정하지 않는다 — 그것이 그림이 제각각으로 보이는 원인이다.
%
% 쓰는 법:  \input{gnc-style}  (figures/src/ 안에서)

\usepackage{tikz}
\usepackage{amsmath}
\usepackage{xcolor}
\usetikzlibrary{arrows.meta, positioning, calc, fit, backgrounds, decorations.pathmorphing}

% ---- 색 --------------------------------------------------------------------
% 색은 강조 하나에만 쓴다. 나머지는 검정.
\definecolor{ink}{HTML}{1A1A1A}
\definecolor{accent}{HTML}{7C3AED}
\definecolor{warn}{HTML}{B91C1C}
\definecolor{soft}{HTML}{666666}

% 기하 그림(좌표계·유도)에서만 쓰는 넷. 블록선도는 위의 흑백 + accent 로 충분하다.
\definecolor{pathc}{HTML}{0D9488}   % 계획 경로 · 항적
\definecolor{errc}{HTML}{C2410C}    % 오차 (cross-track 등)
\definecolor{alng}{HTML}{B45309}    % along-track
\definecolor{flow}{HTML}{0369A1}    % 조류 · 유속

% ---- 선과 화살촉 -----------------------------------------------------------
\tikzset{
  % 신호선. 화살촉은 작고 뾰족한 것 하나뿐이다.
  sig/.style   = {ink, line width=0.5pt, -{Stealth[length=4pt, width=3pt]}},
  wire/.style  = {ink, line width=0.5pt},
  % 강조 경로 — 파선으로 구분한다. 흑백 인쇄에서도 살아야 하기 때문이다.
  asig/.style  = {accent, line width=0.5pt, densely dashed,
                  -{Stealth[length=4pt, width=3pt]}},
  awire/.style = {accent, line width=0.5pt, densely dashed},
  % 블록. 흰 바탕, 직각, 검은 테두리.
  blk/.style   = {draw=ink, line width=0.55pt, fill=white,
                  minimum width=17mm, minimum height=8mm, inner sep=2pt, font=\small},
  % 합산점
  sum/.style   = {draw=ink, line width=0.55pt, fill=white, circle,
                  inner sep=0pt, minimum size=4.6mm, font=\scriptsize},
  asum/.style  = {draw=accent, line width=0.55pt, fill=white, circle,
                  inner sep=0pt, minimum size=4.6mm, font=\scriptsize},
  % 분기점
  dot/.style   = {ink, fill=ink, circle, inner sep=0pt, minimum size=1.5mm},
  % 신호 이름 — 선 위에 작게
  sname/.style = {font=\small, inner sep=1.5pt},
  % 그림 안의 짧은 설명. 그림 안의 글은 최소로 둔다
  note/.style  = {font=\scriptsize, soft, inner sep=1.5pt},
  anote/.style = {font=\scriptsize, accent, inner sep=1.5pt},
  % 축
  axis/.style  = {ink, line width=0.5pt},
  tick/.style  = {font=\scriptsize, soft},
}

% 캡션 한 줄 — 그림 아래 가는 선과 함께
\newcommand{\gnccaption}[2]{%
  \node[anchor=north west, text width=#1, align=left, font=\scriptsize, ink]
        at (current bounding box.south west) {\vspace{2pt}#2};}
  (figures/src/ 안에서)

\usepackage{tikz}
\usepackage{amsmath}
\usepackage{xcolor}
\usetikzlibrary{arrows.meta, positioning, calc, fit, backgrounds, decorations.pathmorphing}

% ---- 색 --------------------------------------------------------------------
% 색은 강조 하나에만 쓴다. 나머지는 검정.
\definecolor{ink}{HTML}{1A1A1A}
\definecolor{accent}{HTML}{7C3AED}
\definecolor{warn}{HTML}{B91C1C}
\definecolor{soft}{HTML}{666666}

% 기하 그림(좌표계·유도)에서만 쓰는 넷. 블록선도는 위의 흑백 + accent 로 충분하다.
\definecolor{pathc}{HTML}{0D9488}   % 계획 경로 · 항적
\definecolor{errc}{HTML}{C2410C}    % 오차 (cross-track 등)
\definecolor{alng}{HTML}{B45309}    % along-track
\definecolor{flow}{HTML}{0369A1}    % 조류 · 유속

% ---- 선과 화살촉 -----------------------------------------------------------
\tikzset{
  % 신호선. 화살촉은 작고 뾰족한 것 하나뿐이다.
  sig/.style   = {ink, line width=0.5pt, -{Stealth[length=4pt, width=3pt]}},
  wire/.style  = {ink, line width=0.5pt},
  % 강조 경로 — 파선으로 구분한다. 흑백 인쇄에서도 살아야 하기 때문이다.
  asig/.style  = {accent, line width=0.5pt, densely dashed,
                  -{Stealth[length=4pt, width=3pt]}},
  awire/.style = {accent, line width=0.5pt, densely dashed},
  % 블록. 흰 바탕, 직각, 검은 테두리.
  blk/.style   = {draw=ink, line width=0.55pt, fill=white,
                  minimum width=17mm, minimum height=8mm, inner sep=2pt, font=\small},
  % 합산점
  sum/.style   = {draw=ink, line width=0.55pt, fill=white, circle,
                  inner sep=0pt, minimum size=4.6mm, font=\scriptsize},
  asum/.style  = {draw=accent, line width=0.55pt, fill=white, circle,
                  inner sep=0pt, minimum size=4.6mm, font=\scriptsize},
  % 분기점
  dot/.style   = {ink, fill=ink, circle, inner sep=0pt, minimum size=1.5mm},
  % 신호 이름 — 선 위에 작게
  sname/.style = {font=\small, inner sep=1.5pt},
  % 그림 안의 짧은 설명. 그림 안의 글은 최소로 둔다
  note/.style  = {font=\scriptsize, soft, inner sep=1.5pt},
  anote/.style = {font=\scriptsize, accent, inner sep=1.5pt},
  % 축
  axis/.style  = {ink, line width=0.5pt},
  tick/.style  = {font=\scriptsize, soft},
}

% 캡션 한 줄 — 그림 아래 가는 선과 함께
\newcommand{\gnccaption}[2]{%
  \node[anchor=north west, text width=#1, align=left, font=\scriptsize, ink]
        at (current bounding box.south west) {\vspace{2pt}#2};}
  (figures/src/ 안에서)

\usepackage{tikz}
\usepackage{amsmath}
\usepackage{xcolor}
\usetikzlibrary{arrows.meta, positioning, calc, fit, backgrounds, decorations.pathmorphing}

% ---- 색 --------------------------------------------------------------------
% 색은 강조 하나에만 쓴다. 나머지는 검정.
\definecolor{ink}{HTML}{1A1A1A}
\definecolor{accent}{HTML}{7C3AED}
\definecolor{warn}{HTML}{B91C1C}
\definecolor{soft}{HTML}{666666}

% 기하 그림(좌표계·유도)에서만 쓰는 넷. 블록선도는 위의 흑백 + accent 로 충분하다.
\definecolor{pathc}{HTML}{0D9488}   % 계획 경로 · 항적
\definecolor{errc}{HTML}{C2410C}    % 오차 (cross-track 등)
\definecolor{alng}{HTML}{B45309}    % along-track
\definecolor{flow}{HTML}{0369A1}    % 조류 · 유속

% ---- 선과 화살촉 -----------------------------------------------------------
\tikzset{
  % 신호선. 화살촉은 작고 뾰족한 것 하나뿐이다.
  sig/.style   = {ink, line width=0.5pt, -{Stealth[length=4pt, width=3pt]}},
  wire/.style  = {ink, line width=0.5pt},
  % 강조 경로 — 파선으로 구분한다. 흑백 인쇄에서도 살아야 하기 때문이다.
  asig/.style  = {accent, line width=0.5pt, densely dashed,
                  -{Stealth[length=4pt, width=3pt]}},
  awire/.style = {accent, line width=0.5pt, densely dashed},
  % 블록. 흰 바탕, 직각, 검은 테두리.
  blk/.style   = {draw=ink, line width=0.55pt, fill=white,
                  minimum width=17mm, minimum height=8mm, inner sep=2pt, font=\small},
  % 합산점
  sum/.style   = {draw=ink, line width=0.55pt, fill=white, circle,
                  inner sep=0pt, minimum size=4.6mm, font=\scriptsize},
  asum/.style  = {draw=accent, line width=0.55pt, fill=white, circle,
                  inner sep=0pt, minimum size=4.6mm, font=\scriptsize},
  % 분기점
  dot/.style   = {ink, fill=ink, circle, inner sep=0pt, minimum size=1.5mm},
  % 신호 이름 — 선 위에 작게
  sname/.style = {font=\small, inner sep=1.5pt},
  % 그림 안의 짧은 설명. 그림 안의 글은 최소로 둔다
  note/.style  = {font=\scriptsize, soft, inner sep=1.5pt},
  anote/.style = {font=\scriptsize, accent, inner sep=1.5pt},
  % 축
  axis/.style  = {ink, line width=0.5pt},
  tick/.style  = {font=\scriptsize, soft},
}

% 캡션 한 줄 — 그림 아래 가는 선과 함께
\newcommand{\gnccaption}[2]{%
  \node[anchor=north west, text width=#1, align=left, font=\scriptsize, ink]
        at (current bounding box.south west) {\vspace{2pt}#2};}


%  2026-10-17: 사용자의 한국어 참고 그림대로 **치수선**을 넣었다. d · R · x_e ·
%  d - x_e 를 길이로 보여 주지 않으면 "남은 거리가 R 보다 작은가" 가 그림에서
%  읽히지 않고 글로만 남는다.
\tikzset{
  dim/.style  = {soft, line width=0.35pt,
                 {Stealth[length=3pt,width=2.2pt]}-{Stealth[length=3pt,width=2.2pt]}},
  ext/.style  = {soft, line width=0.3pt, densely dotted},
  dlbl/.style = {font=\scriptsize, soft, inner sep=1.4pt, fill=white},
}

\begin{document}
\begin{tikzpicture}[x=1mm, y=1mm]

\def\SC{1.5}       % 1 m = 1.5 mm
\def\LEG{78}       % 다리 52 m
\def\RR{7.5}       % R = 5 m
\def\YV{15}        % y_e^{\,p} = 10 m  (= 2R)
\def\XV{71.25}     % d - x_e^{\,p} = 4.5 m  ->  x = 52 - 4.5 = 47.5 m

% ===================== (a) along-track =====================================
\begin{scope}
  \node[font=\small\bfseries, anchor=west] at (0,30)
       {(a)\quad along-track — what MSS uses};

  % 전환 띠 : 남은 거리 < R 인 곳 전부
  \fill[pathc, opacity=0.10] ({\LEG-\RR},-24) rectangle ({\LEG+4},14);
  \draw[pathc, line width=0.4pt, densely dashed] ({\LEG-\RR},-24) -- ({\LEG-\RR},14);
  \node[pathc, font=\scriptsize, anchor=south east] at ({\LEG-\RR-1},4.5)
       {the band $d-x_e^{\,p} < R$};

  %  치수선 — 위에 d 와 R, 아래에 x_e 와 남은 거리
  \draw[ext] (0,0)     -- (0,22);
  \draw[ext] (\LEG,0)  -- (\LEG,22);
  \draw[ext] ({\LEG-\RR},14) -- ({\LEG-\RR},16.5);
  \draw[dim] (0,20) -- (\LEG,20);
  \node[dlbl] at ({0.5*\LEG},20) {$d$ — the leg's length};
  \draw[dim] ({\LEG-\RR},15) -- (\LEG,15);
  \node[dlbl, anchor=south] at ({\LEG-0.5*\RR},15.8) {$R$};

  \draw[ext] (\XV,0) -- (\XV,-10);
  \draw[dim] (0,-5) -- (\XV,-5);
  \node[dlbl] at ({0.5*\XV},-5) {$x_e^{\,p}$ — how far it has come};
  \draw[dim] (\XV,-10) -- (\LEG,-10);
  \node[dlbl, anchor=east] at ({\XV-1.5},-10) {$d-x_e^{\,p}$};

  \draw[pathc, line width=0.9pt, -{Stealth[length=4.5pt,width=3.4pt]}] (0,0) -- (\LEG,0);
  \filldraw[pathc, fill=white, line width=0.7pt] (0,0)     circle (1.1);
  \filldraw[pathc, fill=white, line width=0.7pt] (\LEG,0)  circle (1.1);
  \node[pathc, font=\scriptsize, anchor=north] at (0,-2)    {$\mathbf{p}_i^{\,n}$};
  \node[pathc, font=\scriptsize, anchor=south] at (\LEG,2)  {$\mathbf{p}_{i+1}^{\,n}$};

  \draw[errc, line width=0.8pt] (\XV,0) -- (\XV,-\YV);
  \node[errc, font=\scriptsize, anchor=west] at ({\XV+1.5},{-0.6*\YV}) {$y_e^{\,p} = 2R$};
  \fill[ink] (\XV,{-\YV-1}) -- ({\XV-2.4},{-\YV-6}) -- ({\XV+2.4},{-\YV-6}) -- cycle;

  \node[pathc, font=\small\bfseries, anchor=north] at (\XV,{-\YV-8}) {SWITCHES};

  \node[anchor=north west, align=left, text width=76mm, font=\scriptsize] at (0,-30) {%
    The test looks only at how much of the leg is left, so being far off the path
    cannot defeat it: a vessel blown wide still moves on to the next leg.\\[3pt]
    \texttt{if ( d - x\_e < R\_switch ) k = k + 1}};
\end{scope}

% ===================== (b) circle of acceptance =============================
\begin{scope}[shift={(96,0)}]
  \node[font=\small\bfseries, anchor=west] at (0,30)
       {(b)\quad circle of acceptance — what most books draw};

  \fill[alng, opacity=0.10] (\LEG,0) circle (\RR);
  \draw[alng, line width=0.4pt, densely dashed] (\LEG,0) circle (\RR);
  \draw[alng, line width=0.5pt] (\LEG,0) -- ({\LEG},{\RR});
  \node[alng, font=\scriptsize, anchor=west] at ({\LEG+1},{0.6*\RR}) {$R$};

  \draw[pathc, line width=0.9pt, -{Stealth[length=4.5pt,width=3.4pt]}] (0,0) -- (\LEG,0);
  \filldraw[pathc, fill=white, line width=0.7pt] (0,0)     circle (1.1);
  \filldraw[pathc, fill=white, line width=0.7pt] (\LEG,0)  circle (1.1);
  \node[pathc, font=\scriptsize, anchor=north] at (0,-2)   {$\mathbf{p}_i^{\,n}$};
  \node[pathc, font=\scriptsize, anchor=south] at (\LEG,2) {$\mathbf{p}_{i+1}^{\,n}$};

  %  같은 두 길이를 직각삼각형의 두 변으로 보여 준다 — 빗변이 참 거리다
  \draw[ext] (\XV,-10) -- (\LEG,-10);
  \draw[dim] (\XV,-10) -- (\LEG,-10);
  \node[dlbl, anchor=east] at ({\XV-1.5},-10) {$d-x_e^{\,p}$};
  \draw[errc, line width=0.4pt] (\XV,-2.6) -- ({\XV+2.6},-2.6) -- ({\XV+2.6},0);

  % 참 거리 — 이것이 판정하는 양이다
  \draw[warn, line width=0.7pt, densely dotted] (\XV,-\YV) -- (\LEG,0);
  \node[warn, font=\scriptsize, anchor=north west] at ({\XV+6},{-\YV+5})
       {$\lVert\mathbf{p}-\mathbf{p}_{i+1}^{\,n}\rVert = 10.97$\,m};

  \draw[errc, line width=0.8pt] (\XV,0) -- (\XV,-\YV);
  \fill[ink] (\XV,{-\YV-1}) -- ({\XV-2.4},{-\YV-6}) -- ({\XV+2.4},{-\YV-6}) -- cycle;
  \node[warn, font=\small\bfseries, anchor=north] at (\XV,{-\YV-8}) {DOES NOT SWITCH};

  \node[anchor=north west, align=left, text width=76mm, font=\scriptsize] at (0,-30) {%
    The test is a true distance, so a vessel that never enters the circle waits
    for a waypoint it has already passed — and can wait for ever.\\[3pt]
    \texttt{if ( norm(p - wp\_next) < R ) k = k + 1}};
\end{scope}

% ---- 아래 한 줄 -------------------------------------------------------------
\draw[soft, line width=0.3pt] (0,-46) -- (172,-46);
\node[anchor=north west, align=left, text width=172mm, font=\scriptsize] at (0,-48) {%
  On the path the two agree exactly, because $y_e^{\,p}=0$ makes
  $\lVert\mathbf{p}-\mathbf{p}_{i+1}^{\,n}\rVert = d-x_e^{\,p}$. Off it they do not:
  $\lVert\mathbf{p}-\mathbf{p}_{i+1}^{\,n}\rVert = \sqrt{(d-x_e^{\,p})^2+y_e^2} \ge d-x_e^{\,p}$,
  so \textbf{the circle test is always the stricter of the two}, by an amount that
  grows with $y_e^{\,p}$. Experiment 5-4 evaluates both on one position: the along-track
  remainder is $4.50$\,m and passes; the true distance is $10.97$\,m and fails.
  The one hard limit is $R < \min_k d_k$ — otherwise a leg is already ``reached''
  before the vessel has travelled any of it.};

\end{tikzpicture}
\end{document}
